% Part: sets-functions-relations
% Chapter: functions
% Section: isomorphisms

\documentclass[../../../include/open-logic-section]{subfiles}

\begin{document}

\olfileid{sfr}{fun}{iso}
\olsection{同构}

\begin{explain}
\emph{同构}是保持其所关联集合之结构的双射，而所谓结构，是指集合元素之间成立的关系。考虑下面两个集合 $X=\{1,2,3\}$ 和 $Y=\{4,5,6\}$。这两个集合都由后继、小于和大于关系赋予结构。这两个集合之间的同构，就是保持这些结构的双射。因此，双射函数 $f \colon X \to Y$ 是同构，当且仅当：$i<j$ 当且仅当 $f(i)<f(j)$，$i>j$ 当且仅当 $f(i)>f(j)$，且 $j$ 是 $i$ 的后继当且仅当 $f(j)$ 是 $f(i)$ 的后继。
\end{explain}

\begin{defn}[同构]
设 $U$ 为序对 $\langle X, R\rangle$，$V$ 为序对 $\langle Y, S\rangle$，其中 $X$ 和 $Y$ 是集合，$R$ 和 $S$ 分别是 $X$ 和 $Y$ 上的关系。从 $X$ 到 $Y$ 的双射 $f$ 是从 $U$ 到 $V$ 的\emph{同构}，当且仅当它保持关系结构，即对 $X$ 中的任意 $x_{1}$ 和 $x_{2}$，$\tuple{x_1,x_2} \in R$ 当且仅当 $\tuple{f(x_1),f(x_2)} \in S$。
\end{defn}

\begin{ex}
考虑下面两个集合 $X=\{1,2,3\}$ 和 $Y=\{4,5,6\}$，以及小于和大于关系。函数 $f\colon X \to Y$（其中 $f(x) = 7-x$）是 $\tuple{X,<}$ 与 $\tuple{Y,>}$ 之间的同构。
\end{ex}

\end{document}